\section{Muon Algorithm Background}
\label{sec:muon_background}

Muon~\cite{jordan2024muon} is a spectral optimizer that computes steepest descent updates under the operator (spectral) norm for weight matrices of dimension two or higher.
We summarize its key components.

\paragraph{Steepest descent under spectral norm.}
For a weight matrix $W$ with gradient $G = \nabla_W \mathcal{L}$, the steepest descent direction under the spectral norm is defined as
\begin{equation}
\label{eq:spectral_descent}
\Delta W = \arg\min_{\|\Delta\|_{\mathrm{op}} \leq 1} \langle G, \Delta \rangle,
\end{equation}
where $\|\cdot\|_{\mathrm{op}}$ denotes the operator norm (largest singular value) and $\langle \cdot, \cdot \rangle$ the matrix inner product.
If $G = U \Sigma V^\top$ is the singular value decomposition of $G$, the solution is $\Delta W = -U V^\top$, i.e., the negative sign matrix of the gradient.

\paragraph{Newton--Schulz approximation.}
Computing a full SVD at each step is prohibitively expensive.
Muon instead approximates the orthogonal factor $UV^\top$ via Newton--Schulz iterations for computing the orthogonal polar factor.
The classical iteration $X_{k+1} = \tfrac{1}{2}\, X_k (3I - X_k^\top X_k)$ converges cubically but has a narrow convergence basin.
Muon replaces it with a quintic polynomial variant: starting from
\begin{equation}
\label{eq:ns_init}
X_0 = G / \|G\|_F,
\end{equation}
each iteration applies
\begin{equation}
\label{eq:ns_iteration}
X_{k+1} = a\, X_k + b\, X_k X_k^\top X_k + c\, X_k (X_k^\top X_k)^2,
\end{equation}
with fixed coefficients $a = 3.4445$, $b = -4.7750$, $c = 2.0315$ chosen to maximize the convergence basin.
Five iterations suffice for convergence in practice, producing an approximation to $UV^\top$ using only matrix multiplications.

\paragraph{Momentum and mixed strategy.}
Muon applies Nesterov momentum~\cite{sutskever2013importance} with coefficient $\beta = 0.95$ to the orthogonalized gradient.
The spectral formulation in Eq.~\eqref{eq:spectral_descent} applies only to parameters with matrix structure (2D+); one-dimensional parameters (biases, normalization layer weights and biases, and embedding layers) are handled by a separate AdamW instance with learning rate $3 \times 10^{-4}$ and weight decay $0.05$, yielding a hybrid optimizer.

\paragraph{Relation to Shampoo.}
Both Muon and Shampoo~\cite{gupta2018shampoo} exploit gradient matrix structure, but they differ in mechanism and scope.
Shampoo maintains and periodically inverts Kronecker-factored preconditioners that approximate the full Fisher information matrix, incurring additional memory for the preconditioner state and periodic eigendecomposition steps.
Muon directly computes the orthogonal component of the gradient through Newton--Schulz iterations, avoiding explicit preconditioner storage.
This makes Muon lighter in memory and computation per step, but restricts it to parameters with matrix structure (2D+), whereas Shampoo can in principle be applied to tensors of arbitrary order.
